% !TeX root = ../manual.tex
% ===========================================================================
%  Guía para el administrador
% ===========================================================================

\section{Guía para el administrador}
\label{sec:administrador}

Al iniciar sesión como administrador, el menú lateral muestra siete
secciones, en este orden: \boton{Dashboard}, \boton{Reservas},
\boton{Canchas}, \boton{Disponibilidad}, \boton{Usuarios},
\boton{Configuración} y \boton{Reportes}.

\captura{admin-login}{1}{Inicio de sesión como administrador.}{admin-login}

\subsection{Panel principal (Dashboard)}
\label{sec:admin:dashboard}

Es la primera pantalla al entrar. Resume el día actual con cuatro
indicadores: reservas de hoy, ocupación media de las canchas,
cancelaciones del día y la cancha con mayor demanda hoy.

\captura{admin-dashboard}{1}{Panel principal con los indicadores del día.}{admin-dashboard}

\subsection{Gestión de reservas}
\label{sec:admin:reservas}

Desde \boton{Reservas} se ve el listado completo de reservas del club,
de todos los usuarios, con filtros por fecha, cancha y estado.

\captura{admin-reservas-listado}{1}{Listado global de reservas, con filtros por fecha, cancha y estado.}{admin-reservas-listado}

Cada fila muestra el código de la reserva, el usuario, la cancha, el
deporte, la fecha, el bloque horario y su estado. La columna
\boton{Acciones} solo ofrece \boton{Cancelar} para las reservas en estado
\textsc{Confirmada}.

\subsubsection{Cancelar la reserva de un usuario}

\begin{enumerate}
  \item Localiza la reserva y pulsa \boton{Cancelar} en su fila.
  \item El sistema pide confirmación, mostrando el código, el usuario, la
        cancha y el horario exactos.
\end{enumerate}

\captura{admin-confirmar-cancelacion}{1}{Confirmación antes de cancelar la reserva de un usuario.}{admin-confirmar-cancelacion}

\begin{enumerate}
  \setcounter{enumi}{2}
  \item Pulsa \boton{Cancelar reserva} para confirmar, o \boton{Volver}
        para desistir sin cambiar nada.
\end{enumerate}

La reserva pasa a \textsc{Cancelada} y el bloque horario vuelve a estar
libre de inmediato, tanto en Disponibilidad como en la sesión del usuario
dueño de la reserva.

\captura{admin-reserva-cancelada}{1}{La reserva queda \textsc{Cancelada} tras confirmar.}{admin-reserva-cancelada}

\aviso{Cancelar desde aquí afecta a cualquier usuario del club, no solo a
tus propias reservas. Úsalo cuando la cancelación la pide o la amerita el
usuario (por ejemplo, un reclamo o un cambio de última hora), ya que la
acción es inmediata y no tiene deshacer.}

\subsection{Gestión de canchas}
\label{sec:admin:canchas}

Desde \boton{Canchas} se administra el catálogo completo del club: alta,
edición, horario, bloqueos de mantenimiento y estado activo/inactivo.

\captura{admin-canchas-listado}{1}{Catálogo de canchas, con su horario y estado.}{admin-canchas-listado}

El buscador y los chips de deporte (\boton{Todas}, \boton{Pádel},
\boton{Tenis}, \boton{Básquet}) filtran la lista. Cada tarjeta de cancha
tiene tres acciones, de izquierda a derecha: \boton{Editar}, el icono de
llave inglesa (mantenimiento) y el icono circular (activar/inactivar).

\subsubsection{Crear una cancha}

\begin{enumerate}
  \item Pulsa \boton{+ Nueva Cancha}, arriba a la derecha.
  \item Completa \campo{Nombre}, \campo{Deporte}, y el horario de atención
        (\campo{Apertura} y \campo{Cierre}).
  \item Pulsa \boton{Guardar}.
\end{enumerate}

\captura{admin-nueva-cancha}{1}{Alta de una cancha nueva.}{admin-nueva-cancha}

La cancha aparece de inmediato en el catálogo y queda disponible para
reservar desde ese mismo horario.

\captura{admin-cancha-creada}{1}{La cancha recién creada, ya visible en el catálogo.}{admin-cancha-creada}

\subsubsection{Editar una cancha}

Pulsa \boton{Editar} en la tarjeta de la cancha. Se abre el mismo
formulario que en el alta, con los datos actuales precargados: nombre,
deporte y horario de atención.

\captura{admin-canchas-editar}{1}{Edición de una cancha existente.}{admin-canchas-editar}

\aviso{Cambiar el horario de atención cambia también los bloques que se
ofrecen a partir de ese momento en Disponibilidad. Las reservas ya
confirmadas fuera del nuevo horario no se cancelan automáticamente.}

\subsubsection{Bloquear una cancha por mantenimiento}
\label{sec:admin:mantenimiento}

El icono de llave inglesa abre el panel de \textbf{mantenimiento} de la
cancha: un rango de fechas y horas durante el cual no se admiten reservas
nuevas.

\captura{admin-canchas-mantenimiento}{1}{Panel de mantenimiento: bloqueos vigentes y alta de uno nuevo.}{admin-canchas-mantenimiento}

\begin{enumerate}
  \item Completa \campo{Desde}, \campo{Hasta} y, opcionalmente,
        \campo{Motivo} (por ejemplo, «Mantenimiento de la malla»).
  \item Pulsa \boton{Bloquear}.
\end{enumerate}

El bloqueo aparece en la lista inferior del mismo panel, con la opción
\boton{Retirar} para levantarlo antes de tiempo.

\nota{Un bloqueo de mantenimiento impide reservas \textbf{nuevas} en su
rango; no cancela las reservas que ya estaban confirmadas ahí. Si hace
falta liberar la cancha por completo, cancela también esas reservas desde
Gestión de Reservas (\cref{sec:admin:reservas}).}

\subsubsection{Activar o inactivar una cancha}

El icono circular de la derecha alterna el estado de la cancha:

\begin{itemize}
  \item Con la cancha \textbf{activa}, el icono es una prohibición roja:
        púlsalo para inactivarla.
  \item Con la cancha \textbf{inactiva}, el icono es un visto verde:
        púlsalo para reactivarla.
\end{itemize}

\captura{admin-cancha-inactiva}{1}{Cancha inactiva: deja de ofrecerse en Disponibilidad, sin afectar sus reservas ya hechas.}{admin-cancha-inactiva}

\aviso{Inactivar una cancha la retira de Disponibilidad para nuevas
reservas, pero \textbf{no elimina} las reservas que ya tenía confirmadas
ni su historial. Es la forma correcta de retirar temporalmente una cancha
sin perder información.}

\subsection{Consultar disponibilidad}
\label{sec:admin:disponibilidad}

Desde \boton{Disponibilidad} el administrador ve, para una cancha y una
fecha, exactamente los mismos bloques libres, ocupados y en mantenimiento
que vería un usuario final. Es una pantalla de solo consulta: sirve para
resolver una duda o un reclamo sin tener que reconstruir el listado global
de reservas a mano.

\begin{enumerate}
  \item Entra a \boton{Disponibilidad} desde el menú lateral.
  \item Elige el deporte: \boton{Pádel}, \boton{Tenis} o \boton{Básquet}.
  \item Elige el día, entre hoy y los cuatro días siguientes.
\end{enumerate}

\captura{admin-disponibilidad}{1}{Disponibilidad consultada por el administrador: los mismos bloques que ve un usuario final.}{admin-disponibilidad}

A diferencia de la pantalla del usuario final (\cref{sec:usuario:disponibilidad}),
esta vista tiene dos particularidades:

\begin{itemize}
  \item \textbf{No se puede reservar desde aquí.} Los bloques no son
        pulsables: crear una reserva sigue siendo una acción exclusiva del
        usuario final.
  \item \textbf{También muestra las canchas inactivas}, marcadas con la
        etiqueta \boton{Inactiva} junto al nombre, porque el administrador
        sí las gestiona (\cref{sec:admin:canchas}); al usuario final no se
        le ofrecen.
\end{itemize}

\captura{admin-disponibilidad-inactiva}{1}{Una cancha inactiva aparece marcada como tal, junto a las demás del mismo deporte.}{admin-disponibilidad-inactiva}

\subsection{Gestión de usuarios}
\label{sec:admin:usuarios}

Desde \boton{Usuarios} se ve el listado de todas las cuentas del club,
con su rol y su estado.

\captura{admin-usuarios-listado}{1}{Cuentas del club: rol y estado de cada una.}{admin-usuarios-listado}

La única acción disponible es \boton{Activar} / \boton{Inactivar}, según
el estado actual de la cuenta:

\begin{enumerate}
  \item Localiza la cuenta en la lista.
  \item Pulsa \boton{Inactivar} (cuenta activa) o \boton{Activar} (cuenta
        inactiva).
\end{enumerate}

\captura{admin-usuario-activado}{1}{Activación de una cuenta: pasa a poder iniciar sesión de nuevo.}{admin-usuario-activado}

\nota{Una cuenta inactiva no puede iniciar sesión (véase
\cref{sec:acceso:login}), pero conserva su historial de reservas. Es la
forma de suspender a alguien sin borrar nada de lo que hizo antes.}

\aviso{Esta pantalla no crea cuentas nuevas. Las cuentas se crean desde
\boton{Crear cuenta} en la pantalla de acceso (\cref{sec:acceso:crear-cuenta})
y siempre nacen como usuario final; ascenderlas a administrador se hace
directamente en la base de datos del club, fuera de esta pantalla.}

\subsection{Configuración de reservas activas}
\label{sec:admin:configuracion}

Desde \boton{Configuración} el administrador cambia la política de
reservas del club: el número máximo de reservas confirmadas y futuras
que puede acumular un mismo usuario final al mismo tiempo. Este es el
tope al que se refiere \cref{sec:usuario:reservar}.

\captura{admin-configuracion}{1}{Configuración: el tope vigente de reservas activas por usuario.}{admin-configuracion}

\begin{enumerate}
  \item Escribe el nuevo valor en el campo, junto a \textit{«Tope
        actual»}.
  \item Pulsa \boton{Guardar}. El botón permanece deshabilitado mientras
        el valor no cambie respecto al vigente.
\end{enumerate}

\captura{admin-configuracion-guardado}{1}{Tope actualizado: el aviso confirma el nuevo valor.}{admin-configuracion-guardado}

El cambio rige de inmediato para la siguiente reserva que se intente
crear, sin reiniciar ni redesplegar nada.

\nota{Bajar el tope no cancela nada: los usuarios que ya tuvieran más
reservas activas que el nuevo tope las conservan, y siguen ocupando su
bloque. El tope solo se comprueba al \textbf{crear} una reserva, nunca
sobre las que ya están confirmadas.}

\aviso{El sistema rechaza un tope que no sea un número entero de 1 en
adelante y mantiene vigente el valor anterior. Esta pantalla es exclusiva
del administrador: un usuario final no puede consultarla ni modificarla.}

\subsection{Reportes}
\label{sec:admin:reportes}

Desde \boton{Reportes} se consultan los cuatro indicadores de uso del
club para un rango de fechas: total de reservas, ocupación media,
cancelaciones y la cancha de mayor demanda.

\begin{enumerate}
  \item Ajusta \campo{Desde} y \campo{Hasta}.
  \item Pulsa \boton{Actualizar}.
\end{enumerate}

\captura{admin-reportes-indicadores}{1}{Los cuatro indicadores y el desglose por deporte y por cancha, para el rango elegido.}{admin-reportes-indicadores}

Debajo de los indicadores, dos vistas complementan el análisis:
\textbf{Reservas por deporte} (un conteo por deporte) y
\textbf{Ocupación por cancha} (porcentaje de horas reservadas sobre las
horas de atención de cada cancha en el rango).

Más abajo, el \textbf{Ranking de demanda} ordena todas las canchas por
número de reservas en el rango, señalando además cuál tuvo la menor
demanda.

\captura{admin-reportes-ranking}{1}{Ranking de demanda por cancha, con sus extremos de mayor y menor ocupación.}{admin-reportes-ranking}

\nota{Los reportes se recalculan sobre el rango de fechas elegido cada vez
que pulsas \boton{Actualizar}; no quedan guardados como un informe
aparte.}
